\batchmode
\documentclass[12pt]{article}
\RequirePackage{ifthen}


\usepackage{amsmath}
\usepackage{mathrsfs}
\usepackage{eucal}
\usepackage{amssymb}
\usepackage[french]{babel}
\usepackage[latin1]{inputenc}
\usepackage[T1]{fontenc}




\usepackage[dvips]{color}


\pagecolor[gray]{.7}

\usepackage[latin1]{inputenc}



\makeatletter

\makeatletter
\count@=\the\catcode`\_ \catcode`\_=8 
\newenvironment{tex2html_wrap}{}{}%
\catcode`\<=12\catcode`\_=\count@
\newcommand{\providedcommand}[1]{\expandafter\providecommand\csname #1\endcsname}%
\newcommand{\renewedcommand}[1]{\expandafter\providecommand\csname #1\endcsname{}%
  \expandafter\renewcommand\csname #1\endcsname}%
\newcommand{\newedenvironment}[1]{\newenvironment{#1}{}{}\renewenvironment{#1}}%
\let\newedcommand\renewedcommand
\let\renewedenvironment\newedenvironment
\makeatother
\let\mathon=$
\let\mathoff=$
\ifx\AtBeginDocument\undefined \newcommand{\AtBeginDocument}[1]{}\fi
\newbox\sizebox
\setlength{\hoffset}{0pt}\setlength{\voffset}{0pt}
\addtolength{\textheight}{\footskip}\setlength{\footskip}{0pt}
\addtolength{\textheight}{\topmargin}\setlength{\topmargin}{0pt}
\addtolength{\textheight}{\headheight}\setlength{\headheight}{0pt}
\addtolength{\textheight}{\headsep}\setlength{\headsep}{0pt}
\setlength{\textwidth}{349pt}
\newwrite\lthtmlwrite
\makeatletter
\let\realnormalsize=\normalsize
\global\topskip=2sp
\def\preveqno{}\let\real@float=\@float \let\realend@float=\end@float
\def\@float{\let\@savefreelist\@freelist\real@float}
\def\liih@math{\ifmmode$\else\bad@math\fi}
\def\end@float{\realend@float\global\let\@freelist\@savefreelist}
\let\real@dbflt=\@dbflt \let\end@dblfloat=\end@float
\let\@largefloatcheck=\relax
\let\if@boxedmulticols=\iftrue
\def\@dbflt{\let\@savefreelist\@freelist\real@dbflt}
\def\adjustnormalsize{\def\normalsize{\mathsurround=0pt \realnormalsize
 \parindent=0pt\abovedisplayskip=0pt\belowdisplayskip=0pt}%
 \def\phantompar{\csname par\endcsname}\normalsize}%
\def\lthtmltypeout#1{{\let\protect\string \immediate\write\lthtmlwrite{#1}}}%
\newcommand\lthtmlhboxmathA{\adjustnormalsize\setbox\sizebox=\hbox\bgroup\kern.05em }%
\newcommand\lthtmlhboxmathB{\adjustnormalsize\setbox\sizebox=\hbox to\hsize\bgroup\hfill }%
\newcommand\lthtmlvboxmathA{\adjustnormalsize\setbox\sizebox=\vbox\bgroup %
 \let\ifinner=\iffalse \let\)\liih@math }%
\newcommand\lthtmlboxmathZ{\@next\next\@currlist{}{\def\next{\voidb@x}}%
 \expandafter\box\next\egroup}%
\newcommand\lthtmlmathtype[1]{\gdef\lthtmlmathenv{#1}}%
\newcommand\lthtmllogmath{\dimen0\ht\sizebox \advance\dimen0\dp\sizebox
  \ifdim\dimen0>.95\vsize
   \lthtmltypeout{%
*** image for \lthtmlmathenv\space is too tall at \the\dimen0, reducing to .95 vsize ***}%
   \ht\sizebox.95\vsize \dp\sizebox\z@ \fi
  \lthtmltypeout{l2hSize %
:\lthtmlmathenv:\the\ht\sizebox::\the\dp\sizebox::\the\wd\sizebox.\preveqno}}%
\newcommand\lthtmlfigureA[1]{\let\@savefreelist\@freelist
       \lthtmlmathtype{#1}\lthtmlvboxmathA}%
\newcommand\lthtmlpictureA{\bgroup\catcode`\_=8 \lthtmlpictureB}%
\newcommand\lthtmlpictureB[1]{\lthtmlmathtype{#1}\egroup
       \let\@savefreelist\@freelist \lthtmlhboxmathB}%
\newcommand\lthtmlpictureZ[1]{\hfill\lthtmlfigureZ}%
\newcommand\lthtmlfigureZ{\lthtmlboxmathZ\lthtmllogmath\copy\sizebox
       \global\let\@freelist\@savefreelist}%
\newcommand\lthtmldisplayA{\bgroup\catcode`\_=8 \lthtmldisplayAi}%
\newcommand\lthtmldisplayAi[1]{\lthtmlmathtype{#1}\egroup\lthtmlvboxmathA}%
\newcommand\lthtmldisplayB[1]{\edef\preveqno{(\theequation)}%
  \lthtmldisplayA{#1}\let\@eqnnum\relax}%
\newcommand\lthtmldisplayZ{\lthtmlboxmathZ\lthtmllogmath\lthtmlsetmath}%
\newcommand\lthtmlinlinemathA{\bgroup\catcode`\_=8 \lthtmlinlinemathB}
\newcommand\lthtmlinlinemathB[1]{\lthtmlmathtype{#1}\egroup\lthtmlhboxmathA
  \vrule height1.5ex width0pt }%
\newcommand\lthtmlinlineA{\bgroup\catcode`\_=8 \lthtmlinlineB}%
\newcommand\lthtmlinlineB[1]{\lthtmlmathtype{#1}\egroup\lthtmlhboxmathA}%
\newcommand\lthtmlinlineZ{\egroup\expandafter\ifdim\dp\sizebox>0pt %
  \expandafter\centerinlinemath\fi\lthtmllogmath\lthtmlsetinline}
\newcommand\lthtmlinlinemathZ{\egroup\expandafter\ifdim\dp\sizebox>0pt %
  \expandafter\centerinlinemath\fi\lthtmllogmath\lthtmlsetmath}
\newcommand\lthtmlindisplaymathZ{\egroup %
  \centerinlinemath\lthtmllogmath\lthtmlsetmath}
\def\lthtmlsetinline{\hbox{\vrule width.1em \vtop{\vbox{%
  \kern.1em\copy\sizebox}\ifdim\dp\sizebox>0pt\kern.1em\else\kern.3pt\fi
  \ifdim\hsize>\wd\sizebox \hrule depth1pt\fi}}}
\def\lthtmlsetmath{\hbox{\vrule width.1em\kern-.05em\vtop{\vbox{%
  \kern.1em\kern0.8 pt\hbox{\hglue.17em\copy\sizebox\hglue0.8 pt}}\kern.3pt%
  \ifdim\dp\sizebox>0pt\kern.1em\fi \kern0.8 pt%
  \ifdim\hsize>\wd\sizebox \hrule depth1pt\fi}}}
\def\centerinlinemath{%
  \dimen1=\ifdim\ht\sizebox<\dp\sizebox \dp\sizebox\else\ht\sizebox\fi
  \advance\dimen1by.5pt \vrule width0pt height\dimen1 depth\dimen1 
 \dp\sizebox=\dimen1\ht\sizebox=\dimen1\relax}

\def\lthtmlcheckvsize{\ifdim\ht\sizebox<\vsize 
  \ifdim\wd\sizebox<\hsize\expandafter\hfill\fi \expandafter\vfill
  \else\expandafter\vss\fi}%
\providecommand{\selectlanguage}[1]{}%
\makeatletter \tracingstats = 1 
\providecommand{\Beta}{\textrm{B}}
\providecommand{\Mu}{\textrm{M}}
\providecommand{\Kappa}{\textrm{K}}
\providecommand{\Rho}{\textrm{R}}
\providecommand{\Epsilon}{\textrm{E}}
\providecommand{\Chi}{\textrm{X}}
\providecommand{\Iota}{\textrm{J}}
\providecommand{\omicron}{\textrm{o}}
\providecommand{\Zeta}{\textrm{Z}}
\providecommand{\Eta}{\textrm{H}}
\providecommand{\Nu}{\textrm{N}}
\providecommand{\Omicron}{\textrm{O}}
\providecommand{\Tau}{\textrm{T}}
\providecommand{\Alpha}{\textrm{A}}


\begin{document}
\pagestyle{empty}\thispagestyle{empty}\lthtmltypeout{}%
\lthtmltypeout{latex2htmlLength hsize=\the\hsize}\lthtmltypeout{}%
\lthtmltypeout{latex2htmlLength vsize=\the\vsize}\lthtmltypeout{}%
\lthtmltypeout{latex2htmlLength hoffset=\the\hoffset}\lthtmltypeout{}%
\lthtmltypeout{latex2htmlLength voffset=\the\voffset}\lthtmltypeout{}%
\lthtmltypeout{latex2htmlLength topmargin=\the\topmargin}\lthtmltypeout{}%
\lthtmltypeout{latex2htmlLength topskip=\the\topskip}\lthtmltypeout{}%
\lthtmltypeout{latex2htmlLength headheight=\the\headheight}\lthtmltypeout{}%
\lthtmltypeout{latex2htmlLength headsep=\the\headsep}\lthtmltypeout{}%
\lthtmltypeout{latex2htmlLength parskip=\the\parskip}\lthtmltypeout{}%
\lthtmltypeout{latex2htmlLength oddsidemargin=\the\oddsidemargin}\lthtmltypeout{}%
\makeatletter
\if@twoside\lthtmltypeout{latex2htmlLength evensidemargin=\the\evensidemargin}%
\else\lthtmltypeout{latex2htmlLength evensidemargin=\the\oddsidemargin}\fi%
\lthtmltypeout{}%
\makeatother
\setcounter{page}{1}
\onecolumn

% !!! IMAGES START HERE !!!


%
\providecommand{\card}[1]{\overline{\overline{#1}}}%

{\newpage\clearpage
\lthtmldisplayA{displaymath2468}%
\begin{displaymath}\begin{array}[t]{rcrcl}
      f & : & \left[0,1\right] & \to     & \left]0,1\right] \\\\
        &   & x                & \mapsto & \left\lbrace
                                           \begin{array}{ll}
                                           1/2              & \textrm{si $x=0$ ;} \\
                                           \frac{1}{n+1}    & \textrm{si $x=\frac{1}{n}$ pour un $n > 1$ ;} \\
                                           x                & \textrm{sinon}
                                           \end{array}
                                           \right.
      \end{array}\end{displaymath}%
\lthtmldisplayZ
\lthtmlcheckvsize\clearpage}

{\newpage\clearpage
\lthtmldisplayA{displaymath2470}%
\begin{displaymath}\begin{array}[t]{rcrcl}
      f & : & \left[0,1\right] & \to     & \left]0,1\right[ \\\\
        &   & x                & \mapsto & \left\lbrace
                                           \begin{array}{ll}
                                           1/2              & \textrm{si $x=0$ ;} \\
                                           1/3              & \textrm{si $x=1$ ;} \\
                                           \frac{1}{n+2}    & \textrm{si $x=\frac{1}{n}$ pour un $n > 1$ ;} \\
                                           x                & \textrm{sinon}
                                           \end{array}
                                           \right.
      \end{array}\end{displaymath}%
\lthtmldisplayZ
\lthtmlcheckvsize\clearpage}

{\newpage\clearpage
\lthtmldisplayA{displaymath2472}%
\begin{displaymath}\begin{array}[t]{rcrcl}
      f & : & \mathbb{Q} & \to     & \mathbb{Q}^*_+ \\\\
        &   & x          & \mapsto & \left\lbrace
                                     \begin{array}{ll}
                                     \frac{1}{1-x}          & \textrm{si $x \leqslant 0$ ;} \\
                                     x+1                    & \textrm{si $x>0$}
                                     \end{array}
                                     \right.
      \end{array}\end{displaymath}%
\lthtmldisplayZ
\lthtmlcheckvsize\clearpage}

{\newpage\clearpage
\lthtmlinlinemathA{tex2html_wrap_inline2474}%
$ I_o = \{ \left]a,b \right[ \subseteq \mathbb{R} \mid a,b \in \mathbb{Q} \}$%
\lthtmlinlinemathZ
\lthtmlcheckvsize\clearpage}

{\newpage\clearpage
\lthtmldisplayA{displaymath2476}%
\begin{displaymath}\begin{array}[t]{rcrcl}
      f & : & \mathbb{Q} & \to     & I_o \\
        &   & x          & \mapsto & \left]x,x+1\right[
      \end{array}\end{displaymath}%
\lthtmldisplayZ
\lthtmlcheckvsize\clearpage}

{\newpage\clearpage
\lthtmlinlinemathA{tex2html_wrap_inline2478}%
$ f$%
\lthtmlinlinemathZ
\lthtmlcheckvsize\clearpage}

{\newpage\clearpage
\lthtmlinlinemathA{tex2html_wrap_inline2480}%
$ \aleph_0 = \overline{\overline{\mathbb{Q}}} \leqslant \overline{\overline{I_o}}$%
\lthtmlinlinemathZ
\lthtmlcheckvsize\clearpage}

{\newpage\clearpage
\lthtmldisplayA{displaymath2482}%
\begin{displaymath}\begin{array}[t]{rcrcl}
       g & : & I_o               & \to     & \mathbb{Q} \times \mathbb{Q} \\
         &   & \left]a,b\right[  & \mapsto & (a,b)
        \end{array}\end{displaymath}%
\lthtmldisplayZ
\lthtmlcheckvsize\clearpage}

{\newpage\clearpage
\lthtmlinlinemathA{tex2html_wrap_inline2484}%
$ g$%
\lthtmlinlinemathZ
\lthtmlcheckvsize\clearpage}

{\newpage\clearpage
\lthtmlinlinemathA{tex2html_wrap_inline2486}%
$ \overline{\overline{I_o}} \leqslant \overline{\overline{\mathbb{Q} \times \mathbb{Q}}} = \aleph_0$%
\lthtmlinlinemathZ
\lthtmlcheckvsize\clearpage}

{\newpage\clearpage
\lthtmlinlinemathA{tex2html_wrap_inline2488}%
$ \overline{\overline{I_0}} = \aleph_0$%
\lthtmlinlinemathZ
\lthtmlcheckvsize\clearpage}

{\newpage\clearpage
\lthtmlinlinemathA{tex2html_wrap_inline2498}%
$ \mathbb{N}^n$%
\lthtmlinlinemathZ
\lthtmlcheckvsize\clearpage}

{\newpage\clearpage
\lthtmlinlinemathA{tex2html_wrap_inline2500}%
$ n \in \mathbb{N}^*$%
\lthtmlinlinemathZ
\lthtmlcheckvsize\clearpage}

{\newpage\clearpage
\lthtmlinlinemathA{tex2html_wrap_inline2502}%
$ n=1$%
\lthtmlinlinemathZ
\lthtmlcheckvsize\clearpage}

{\newpage\clearpage
\lthtmlinlinemathA{tex2html_wrap_inline2506}%
$ \mathbb{N}^{n+1}$%
\lthtmlinlinemathZ
\lthtmlcheckvsize\clearpage}

{\newpage\clearpage
\lthtmlinlinemathA{tex2html_wrap_inline2508}%
$ \mathbb{N}^{n+1} = \mathbb{N}^n \times \mathbb{N}$%
\lthtmlinlinemathZ
\lthtmlcheckvsize\clearpage}

{\newpage\clearpage
\lthtmlinlinemathA{tex2html_wrap_inline2512}%
$ \mathbb{N}$%
\lthtmlinlinemathZ
\lthtmlcheckvsize\clearpage}

{\newpage\clearpage
\lthtmlinlinemathA{tex2html_wrap_inline2520}%
$ \mathbb{Z}[x]$%
\lthtmlinlinemathZ
\lthtmlcheckvsize\clearpage}

{\newpage\clearpage
\lthtmlinlinemathA{tex2html_wrap_inline2522}%
$ \mathbb{Z}_n[x] = \{ p \in \mathbb{Z}[x] \mid \textrm{degr\'e}(p) \leqslant n
\}$%
\lthtmlinlinemathZ
\lthtmlcheckvsize\clearpage}

{\newpage\clearpage
\lthtmldisplayA{displaymath2526}%
\begin{displaymath}\begin{array}[t]{rcl}
      \mathbb{Z}_n[x]                  & \to     & \mathbb{Z}^{n+1} \\
      a_0+a_1x+a_2x^2+\ldots+a_nx^n    & \mapsto & (a_0,a_1,a_2,\ldots,a_n)
      \end{array}\end{displaymath}%
\lthtmldisplayZ
\lthtmlcheckvsize\clearpage}

{\newpage\clearpage
\lthtmlinlinemathA{tex2html_wrap_inline2530}%
$ \overline{\overline{\mathbb{Z}_n[x]}} =
\overline{\overline{\mathbb{Z}^{n+1}}}$%
\lthtmlinlinemathZ
\lthtmlcheckvsize\clearpage}

{\newpage\clearpage
\lthtmlinlinemathA{tex2html_wrap_inline2532}%
$ \mathbb{Z}^{n+1}$%
\lthtmlinlinemathZ
\lthtmlcheckvsize\clearpage}

{\newpage\clearpage
\lthtmlinlinemathA{tex2html_wrap_inline2534}%
$ \mathbb{Z}_n[x]$%
\lthtmlinlinemathZ
\lthtmlcheckvsize\clearpage}

{\newpage\clearpage
\lthtmlinlinemathA{tex2html_wrap_inline2536}%
$ \mathbb{Z}[x] = \bigcup\limits_{n \in \mathbb{N}}
\mathbb{Z}_n[x]$%
\lthtmlinlinemathZ
\lthtmlcheckvsize\clearpage}

{\newpage\clearpage
\lthtmlinlinemathA{tex2html_wrap_inline2540}%
$ {\cal A}$%
\lthtmlinlinemathZ
\lthtmlcheckvsize\clearpage}

{\newpage\clearpage
\lthtmlinlinemathA{tex2html_wrap_inline2542}%
$ \mathcal{R}_n$%
\lthtmlinlinemathZ
\lthtmlcheckvsize\clearpage}

{\newpage\clearpage
\lthtmlinlinemathA{tex2html_wrap_inline2548}%
$ n$%
\lthtmlinlinemathZ
\lthtmlcheckvsize\clearpage}

{\newpage\clearpage
\lthtmlinlinemathA{tex2html_wrap_inline2552}%
$ \mathcal{A} = \bigcup\limits_{n \in \mathbb{N}}
\mathcal{R}_n$%
\lthtmlinlinemathZ
\lthtmlcheckvsize\clearpage}

{\newpage\clearpage
\lthtmlinlinemathA{tex2html_wrap_inline2554}%
$ \mathcal{A}$%
\lthtmlinlinemathZ
\lthtmlcheckvsize\clearpage}

{\newpage\clearpage
\lthtmlinlinemathA{tex2html_wrap_inline2556}%
$ {\cal I}$%
\lthtmlinlinemathZ
\lthtmlcheckvsize\clearpage}

{\newpage\clearpage
\lthtmlinlinemathA{tex2html_wrap_inline2558}%
$ j$%
\lthtmlinlinemathZ
\lthtmlcheckvsize\clearpage}

{\newpage\clearpage
\lthtmlinlinemathA{tex2html_wrap_inline2562}%
$ \overline{\overline{\mathcal{I}}}\leqslant \overline{\overline{\mathbb{R}}} = c$%
\lthtmlinlinemathZ
\lthtmlcheckvsize\clearpage}

{\newpage\clearpage
\lthtmlinlinemathA{tex2html_wrap_inline2564}%
$ \mathcal{I}$%
\lthtmlinlinemathZ
\lthtmlcheckvsize\clearpage}

{\newpage\clearpage
\lthtmlinlinemathA{tex2html_wrap_inline2566}%
$ \mathbb{R} = \mathcal{I} \cup \mathbb{Q}$%
\lthtmlinlinemathZ
\lthtmlcheckvsize\clearpage}

{\newpage\clearpage
\lthtmlinlinemathA{tex2html_wrap_inline2570}%
$ {\cal T}$%
\lthtmlinlinemathZ
\lthtmlcheckvsize\clearpage}

{\newpage\clearpage
\lthtmlinlinemathA{tex2html_wrap_inline2576}%
$ \overline{\overline{\mathcal{T}}}\leqslant \overline{\overline{\mathbb{R}}} = c$%
\lthtmlinlinemathZ
\lthtmlcheckvsize\clearpage}

{\newpage\clearpage
\lthtmlinlinemathA{tex2html_wrap_inline2578}%
$ \mathcal{T}$%
\lthtmlinlinemathZ
\lthtmlcheckvsize\clearpage}

{\newpage\clearpage
\lthtmlinlinemathA{tex2html_wrap_inline2580}%
$ \mathbb{R} = \mathcal{T} \cup \mathbb{A}$%
\lthtmlinlinemathZ
\lthtmlcheckvsize\clearpage}

{\newpage\clearpage
\lthtmlinlinemathA{tex2html_wrap_inline2584}%
$ \mathbb{N}^{[n]}$%
\lthtmlinlinemathZ
\lthtmlcheckvsize\clearpage}

{\newpage\clearpage
\lthtmldisplayA{displaymath2586}%
\begin{displaymath}\begin{array}[t]{rcrcl}
          f & : & \mathbb{N} & \to     & \mathbb{N}^{[n]} \\
            &   & n          & \mapsto & \begin{array}[t]{rcrcl}
                                         f_n & : & [n] & \to     & \mathbb{N} \\
                                             &   &  m  & \mapsto & n.
                                         \end{array}
          \end{array}\end{displaymath}%
\lthtmldisplayZ
\lthtmlcheckvsize\clearpage}

{\newpage\clearpage
\lthtmlinlinemathA{tex2html_wrap_inline2590}%
$ \overline{\overline{\mathbb{N}}} \leqslant \overline{\overline{\mathbb{N}^{[n]}}}$%
\lthtmlinlinemathZ
\lthtmlcheckvsize\clearpage}

{\newpage\clearpage
\lthtmldisplayA{displaymath2592}%
\begin{displaymath}\begin{array}[t]{rcrcl}
                 g & : & \mathbb{N}^{[n]} & \to     & \mathbb{N} \\
                   &   &  h               & \mapsto & 2^{h(1)}\cdot 3^{h(2)}\cdot 5^{h(3)}\cdot 7^{h(4)}\cdot 11^{h(5)}\cdot \ldots\cdot  p_n^{h(n)}
          \end{array}\end{displaymath}%
\lthtmldisplayZ
\lthtmlcheckvsize\clearpage}

{\newpage\clearpage
\lthtmlinlinemathA{tex2html_wrap_inline2594}%
$ p_n$%
\lthtmlinlinemathZ
\lthtmlcheckvsize\clearpage}

{\newpage\clearpage
\lthtmlinlinemathA{tex2html_wrap_inline2600}%
$ \overline{\overline{\mathbb{N}^{[n]}}} \leqslant \overline{\overline{\mathbb{N}}}$%
\lthtmlinlinemathZ
\lthtmlcheckvsize\clearpage}

{\newpage\clearpage
\lthtmlinlinemathA{tex2html_wrap_inline2604}%
$ {\cal P}_f(\mathbb{N})$%
\lthtmlinlinemathZ
\lthtmlcheckvsize\clearpage}

{\newpage\clearpage
\lthtmlinlinemathA{tex2html_wrap_inline2606}%
$ \mathcal{P}_n(\mathbb{N})$%
\lthtmlinlinemathZ
\lthtmlcheckvsize\clearpage}

{\newpage\clearpage
\lthtmldisplayA{displaymath2612}%
\begin{displaymath}\begin{array}[t]{rcrcl}
           g & : & \mathbb{N}^{n}         & \to     & \mathcal{P}_n(\mathbb{N}) \\
             &   & (x_1,x_2,\ldots ,x_n)  & \mapsto & \{ x_1, x_2, \ldots, x_n \}
          \end{array}\end{displaymath}%
\lthtmldisplayZ
\lthtmlcheckvsize\clearpage}

{\newpage\clearpage
\lthtmlinlinemathA{tex2html_wrap_inline2616}%
$ \overline{\overline{\mathbb{N}^{n}}} \geqslant \overline{\overline{\mathcal{P}_n(\mathbb{N})}}$%
\lthtmlinlinemathZ
\lthtmlcheckvsize\clearpage}

{\newpage\clearpage
\lthtmlinlinemathA{tex2html_wrap_inline2622}%
$ {\cal P}_f(\mathbb{N}) = \bigcup\limits_{n \in \mathbb{N}} \mathcal{P}_n(\mathbb{N})$%
\lthtmlinlinemathZ
\lthtmlcheckvsize\clearpage}

{\newpage\clearpage
\lthtmldisplayA{displaymath2626}%
\begin{displaymath}\begin{array}[t]{rcrcl}
                  h & : & \mathbb{N} & \to     & \mathcal{P}_f(\mathbb{N}) \\
                  &   & n            & \mapsto & \{ n \}
                  \end{array}\end{displaymath}%
\lthtmldisplayZ
\lthtmlcheckvsize\clearpage}

{\newpage\clearpage
\lthtmlinlinemathA{tex2html_wrap_inline2628}%
$ h$%
\lthtmlinlinemathZ
\lthtmlcheckvsize\clearpage}

{\newpage\clearpage
\lthtmlinlinemathA{tex2html_wrap_inline2630}%
$ \overline{\overline{\mathbb{N}}} \leqslant \overline{\overline{\mathcal{P}_f(\mathbb{N})}}$%
\lthtmlinlinemathZ
\lthtmlcheckvsize\clearpage}

{\newpage\clearpage
\lthtmlinlinemathA{tex2html_wrap_inline2632}%
$ \mathcal{P}_f(\mathbb{N})$%
\lthtmlinlinemathZ
\lthtmlcheckvsize\clearpage}

{\newpage\clearpage
\lthtmlinlinemathA{tex2html_wrap_inline2634}%
$ \mathbb{N}^{\mathbb{N}}$%
\lthtmlinlinemathZ
\lthtmlcheckvsize\clearpage}

{\newpage\clearpage
\lthtmldisplayA{displaymath2636}%
\begin{displaymath}\begin{array}[t]{rcrcl}
          h & : & \mathbb{N}^{\mathbb{N}} & \to     & \mathcal{P}(\mathbb{N}) \setminus \varnothing \\
            &   & f                       & \mapsto & {\rm Im}(f)
          \end{array}\end{displaymath}%
\lthtmldisplayZ
\lthtmlcheckvsize\clearpage}

{\newpage\clearpage
\lthtmlinlinemathA{tex2html_wrap_inline2640}%
$ \overline{\overline{\mathbb{N}^{\mathbb{N}}}} \geqslant \overline{\overline{\mathcal{P}(\mathbb{N}) \setminus \varnothing}}$%
\lthtmlinlinemathZ
\lthtmlcheckvsize\clearpage}

{\newpage\clearpage
\lthtmlinlinemathA{tex2html_wrap_inline2642}%
$ \overline{\overline{\mathcal{P}(\mathbb{N}) \setminus \varnothing}} = \overline{\overline{\mathcal{P}(\mathbb{N})}}
> \overline{\overline{\mathbb{N}}}$%
\lthtmlinlinemathZ
\lthtmlcheckvsize\clearpage}

{\newpage\clearpage
\lthtmlinlinemathA{tex2html_wrap_inline2646}%
$ x \in \mathbb{R}$%
\lthtmlinlinemathZ
\lthtmlcheckvsize\clearpage}

{\newpage\clearpage
\lthtmldisplayA{displaymath2648}%
\begin{displaymath}\begin{array}{rcl}
\widetilde{x}   & = & \{ y \in \mathbb{R} \mid y-x \in \mathbb{Z} \} \\
                & = & \{ y \in \mathbb{R} \mid y-x = k \textrm{ pour un } k \in \mathbb{Z} \} \\
                & = & \{ y \in \mathbb{R} \mid y = x + k \textrm{ pour un } k \in \mathbb{Z} \} \\
                & = & \{ x+k \mid k \in \mathbb{Z} \}
\end{array}\end{displaymath}%
\lthtmldisplayZ
\lthtmlcheckvsize\clearpage}

{\newpage\clearpage
\lthtmlinlinemathA{tex2html_wrap_inline2650}%
$ \overline{\overline{\widetilde{x}}} = \overline{\overline{\mathbb{Z}}} = \aleph_0$%
\lthtmlinlinemathZ
\lthtmlcheckvsize\clearpage}

{\newpage\clearpage
\lthtmlinlinemathA{tex2html_wrap_inline2652}%
$ \left[0,1\right[$%
\lthtmlinlinemathZ
\lthtmlcheckvsize\clearpage}

{\newpage\clearpage
\lthtmlinlinemathA{tex2html_wrap_inline2654}%
$ a \in \mathbb{R}$%
\lthtmlinlinemathZ
\lthtmlcheckvsize\clearpage}

{\newpage\clearpage
\lthtmlinlinemathA{tex2html_wrap_inline2656}%
$ \lfloor a \rfloor$%
\lthtmlinlinemathZ
\lthtmlcheckvsize\clearpage}

{\newpage\clearpage
\lthtmlinlinemathA{tex2html_wrap_inline2658}%
$ a$%
\lthtmlinlinemathZ
\lthtmlcheckvsize\clearpage}

{\newpage\clearpage
\lthtmlinlinemathA{tex2html_wrap_inline2660}%
$ a - \lfloor a \rfloor \in
\left[0,1\right[$%
\lthtmlinlinemathZ
\lthtmlcheckvsize\clearpage}

{\newpage\clearpage
\lthtmlinlinemathA{tex2html_wrap_inline2662}%
$ (a - \lfloor a \rfloor) - a = \lfloor
a \rfloor \in \mathbb{Z}$%
\lthtmlinlinemathZ
\lthtmlcheckvsize\clearpage}

{\newpage\clearpage
\lthtmlinlinemathA{tex2html_wrap_inline2664}%
$ a - \lfloor a \rfloor \in
\widetilde{a} \cap \left[0,1\right[$%
\lthtmlinlinemathZ
\lthtmlcheckvsize\clearpage}

{\newpage\clearpage
\lthtmlinlinemathA{tex2html_wrap_inline2666}%
$ \widetilde{a} \cap \left[0,1\right[$%
\lthtmlinlinemathZ
\lthtmlcheckvsize\clearpage}

{\newpage\clearpage
\lthtmlinlinemathA{tex2html_wrap_inline2668}%
$ u,v \in
\widetilde{a} \cap \left[0,1\right[$%
\lthtmlinlinemathZ
\lthtmlcheckvsize\clearpage}

{\newpage\clearpage
\lthtmlinlinemathA{tex2html_wrap_inline2670}%
$ u \in \widetilde{a}
\Rightarrow u-a \in \mathbb{Z}$%
\lthtmlinlinemathZ
\lthtmlcheckvsize\clearpage}

{\newpage\clearpage
\lthtmlinlinemathA{tex2html_wrap_inline2672}%
$ v \in \widetilde{a}
\Rightarrow a-v \in \mathbb{Z}$%
\lthtmlinlinemathZ
\lthtmlcheckvsize\clearpage}

{\newpage\clearpage
\lthtmlinlinemathA{tex2html_wrap_inline2674}%
$ (u-a)+(a-v) = u-v
\in \mathbb{Z}$%
\lthtmlinlinemathZ
\lthtmlcheckvsize\clearpage}

{\newpage\clearpage
\lthtmlinlinemathA{tex2html_wrap_inline2676}%
$ u,v \in \left[0,1\right[$%
\lthtmlinlinemathZ
\lthtmlcheckvsize\clearpage}

{\newpage\clearpage
\lthtmlinlinemathA{tex2html_wrap_inline2678}%
$ u-v \in
\mathbb{Z}$%
\lthtmlinlinemathZ
\lthtmlcheckvsize\clearpage}

{\newpage\clearpage
\lthtmlinlinemathA{tex2html_wrap_inline2680}%
$ u=v$%
\lthtmlinlinemathZ
\lthtmlcheckvsize\clearpage}

{\newpage\clearpage
\lthtmldisplayA{displaymath2682}%
\begin{displaymath}\begin{array}[t]{rcrcl}
          h & : & \mathbb{R}/S  & \to     & \left[0,1\right[ \\
            &   & \widetilde{a} & \mapsto & a - \lfloor a \rfloor
          \end{array}\end{displaymath}%
\lthtmldisplayZ
\lthtmlcheckvsize\clearpage}

{\newpage\clearpage
\lthtmlinlinemathA{tex2html_wrap_inline2686}%
$ \overline{\overline{\mathbb{R}/S}} =
\overline{\overline{\left[0,1\right[}} = c$%
\lthtmlinlinemathZ
\lthtmlcheckvsize\clearpage}

{\newpage\clearpage
\lthtmlinlinemathA{tex2html_wrap_inline2690}%
$ i \in I$%
\lthtmlinlinemathZ
\lthtmlcheckvsize\clearpage}

{\newpage\clearpage
\lthtmlinlinemathA{tex2html_wrap_inline2692}%
$ A_i = \left]a_i,b_i\right[$%
\lthtmlinlinemathZ
\lthtmlcheckvsize\clearpage}

{\newpage\clearpage
\lthtmlinlinemathA{tex2html_wrap_inline2694}%
$ a_i,b_i \in \mathbb{R}$%
\lthtmlinlinemathZ
\lthtmlcheckvsize\clearpage}

{\newpage\clearpage
\lthtmlinlinemathA{tex2html_wrap_inline2696}%
$ q_i \in \mathbb{Q}$%
\lthtmlinlinemathZ
\lthtmlcheckvsize\clearpage}

{\newpage\clearpage
\lthtmlinlinemathA{tex2html_wrap_inline2698}%
$ a_i$%
\lthtmlinlinemathZ
\lthtmlcheckvsize\clearpage}

{\newpage\clearpage
\lthtmlinlinemathA{tex2html_wrap_inline2700}%
$ b_i$%
\lthtmlinlinemathZ
\lthtmlcheckvsize\clearpage}

{\newpage\clearpage
\lthtmldisplayA{displaymath2702}%
\begin{displaymath}\begin{array}[t]{rcrcl}
                         f & : & I  & \to     & \mathbb{Q} \\
                     &   & i  & \mapsto & q_i
          \end{array}\end{displaymath}%
\lthtmldisplayZ
\lthtmlcheckvsize\clearpage}

{\newpage\clearpage
\lthtmlinlinemathA{tex2html_wrap_inline2706}%
$ \overline{\overline{I}} \leqslant \overline{\overline{\mathbb{Q}}} = \aleph_0$%
\lthtmlinlinemathZ
\lthtmlcheckvsize\clearpage}

{\newpage\clearpage
\lthtmlinlinemathA{tex2html_wrap_inline2708}%
$ I$%
\lthtmlinlinemathZ
\lthtmlcheckvsize\clearpage}

{\newpage\clearpage
\lthtmlinlinemathA{tex2html_wrap_inline2712}%
$ D_i = \{ (x,y) \in \mathbb{R}^2 \mid (x-a_i)^2 + (y-b_i)^2 < r_i^2$%
\lthtmlinlinemathZ
\lthtmlcheckvsize\clearpage}

{\newpage\clearpage
\lthtmlinlinemathA{tex2html_wrap_inline2714}%
$ a_i,b_i$%
\lthtmlinlinemathZ
\lthtmlcheckvsize\clearpage}

{\newpage\clearpage
\lthtmlinlinemathA{tex2html_wrap_inline2716}%
$ r_i \in \mathbb{R}$%
\lthtmlinlinemathZ
\lthtmlcheckvsize\clearpage}

{\newpage\clearpage
\lthtmlinlinemathA{tex2html_wrap_inline2718}%
$ (p_i,q_i) \in \mathbb{Q}^2$%
\lthtmlinlinemathZ
\lthtmlcheckvsize\clearpage}

{\newpage\clearpage
\lthtmlinlinemathA{tex2html_wrap_inline2720}%
$ D_i$%
\lthtmlinlinemathZ
\lthtmlcheckvsize\clearpage}

{\newpage\clearpage
\lthtmldisplayA{displaymath2722}%
\begin{displaymath}\begin{array}[t]{rcrcl}
                         g & : & I  & \to     & \mathbb{Q}^2 \\
                     &   & i    & \mapsto & (p_i,q_i)
          \end{array}\end{displaymath}%
\lthtmldisplayZ
\lthtmlcheckvsize\clearpage}

{\newpage\clearpage
\lthtmlinlinemathA{tex2html_wrap_inline2726}%
$ \overline{\overline{I}} \leqslant \overline{\overline{\mathbb{Q}^2}} = \aleph_0$%
\lthtmlinlinemathZ
\lthtmlcheckvsize\clearpage}

{\newpage\clearpage
\lthtmlinlinemathA{tex2html_wrap_inline2730}%
$ \mathcal{S}_f$%
\lthtmlinlinemathZ
\lthtmlcheckvsize\clearpage}

{\newpage\clearpage
\lthtmlinlinemathA{tex2html_wrap_inline2732}%
$ \mathcal{S}_n$%
\lthtmlinlinemathZ
\lthtmlcheckvsize\clearpage}

{\newpage\clearpage
\lthtmlinlinemathA{tex2html_wrap_inline2736}%
$ \overline{\overline{\mathcal{S}_n}} = \overline{\overline{\prod\limits_{i=1}^{n} \mathbb{N}}} = \aleph_0$%
\lthtmlinlinemathZ
\lthtmlcheckvsize\clearpage}

{\newpage\clearpage
\lthtmlinlinemathA{tex2html_wrap_inline2738}%
$ \mathcal{S}_f = \bigcup\limits_{n \in \mathbb{N}} \mathcal{S}_n$%
\lthtmlinlinemathZ
\lthtmlcheckvsize\clearpage}

{\newpage\clearpage
\lthtmldisplayA{displaymath2750}%
\begin{displaymath}\begin{array}[t]{rcrcl}
             f & : & \mathbb{R}^+_*  & \to     & \mathcal{I} \\
               &   & x                     & \mapsto & \begin{array}[t]{ll}
                                                                                 x          & \textrm{si $x$\  est irrationnel} \\
                                                                                 -x\sqrt{2} & \textrm{si $x$\  est rationnel}
                                                                                 \end{array}
             \end{array}\end{displaymath}%
\lthtmldisplayZ
\lthtmlcheckvsize\clearpage}

{\newpage\clearpage
\lthtmlinlinemathA{tex2html_wrap_inline2754}%
$ \overline{\overline{\mathbb{R}^+_*}} \leqslant \overline{\overline{\mathcal{I}}}$%
\lthtmlinlinemathZ
\lthtmlcheckvsize\clearpage}

{\newpage\clearpage
\lthtmlinlinemathA{tex2html_wrap_inline2756}%
$ c \leqslant \overline{\overline{\mathcal{I}}}$%
\lthtmlinlinemathZ
\lthtmlcheckvsize\clearpage}

{\newpage\clearpage
\lthtmlinlinemathA{tex2html_wrap_inline2758}%
$ \overline{\overline{\mathcal{I}}} = c$%
\lthtmlinlinemathZ
\lthtmlcheckvsize\clearpage}

{\newpage\clearpage
\lthtmldisplayA{displaymath2760}%
\begin{displaymath}\begin{array}[t]{rcrcl}
                         f & : & \mathcal{B}  & \to     & \mathbb{R} \times \mathbb{R} \times \mathbb{R}^+ \\
                     &   & D                    & \mapsto & (x_D, y_D, r_D)
          \end{array}\end{displaymath}%
\lthtmldisplayZ
\lthtmlcheckvsize\clearpage}

{\newpage\clearpage
\lthtmlinlinemathA{tex2html_wrap_inline2762}%
$ (x_D, y_D)$%
\lthtmlinlinemathZ
\lthtmlcheckvsize\clearpage}

{\newpage\clearpage
\lthtmlinlinemathA{tex2html_wrap_inline2764}%
$ D$%
\lthtmlinlinemathZ
\lthtmlcheckvsize\clearpage}

{\newpage\clearpage
\lthtmlinlinemathA{tex2html_wrap_inline2766}%
$ r_D$%
\lthtmlinlinemathZ
\lthtmlcheckvsize\clearpage}

{\newpage\clearpage
\lthtmlinlinemathA{tex2html_wrap_inline2770}%
$ \overline{\overline{\mathcal{B}}} = \overline{\overline{\mathbb{R} \times \mathbb{R} \times \mathbb{R}^+}} = \overline{\overline{\mathbb{R}}} = c$%
\lthtmlinlinemathZ
\lthtmlcheckvsize\clearpage}

{\newpage\clearpage
\lthtmldisplayA{displaymath2772}%
\begin{displaymath}\begin{array}[t]{rcrcl}
                    f & : & \mathbb{R}^2  & \to     & \mathcal{T} \\
                      &   & (x,y)         & \mapsto & T_{(x,y)}
          \end{array}\end{displaymath}%
\lthtmldisplayZ
\lthtmlcheckvsize\clearpage}

{\newpage\clearpage
\lthtmlinlinemathA{tex2html_wrap_inline2774}%
$ T_{(x,y)}$%
\lthtmlinlinemathZ
\lthtmlcheckvsize\clearpage}

{\newpage\clearpage
\lthtmlinlinemathA{tex2html_wrap_inline2776}%
$ (x,y)$%
\lthtmlinlinemathZ
\lthtmlcheckvsize\clearpage}

{\newpage\clearpage
\lthtmlinlinemathA{tex2html_wrap_inline2778}%
$ (x,y+1)$%
\lthtmlinlinemathZ
\lthtmlcheckvsize\clearpage}

{\newpage\clearpage
\lthtmlinlinemathA{tex2html_wrap_inline2780}%
$ (x+1,y)$%
\lthtmlinlinemathZ
\lthtmlcheckvsize\clearpage}

{\newpage\clearpage
\lthtmlinlinemathA{tex2html_wrap_inline2784}%
$ c = \overline{\overline{\mathbb{R}^2}} \leqslant
\overline{\overline{\mathcal{T}}}$%
\lthtmlinlinemathZ
\lthtmlcheckvsize\clearpage}

{\newpage\clearpage
\lthtmlinlinemathA{tex2html_wrap_inline2786}%
$ \mathbb{R}^2 \times \mathbb{R}^2 \times
\mathbb{R}^2$%
\lthtmlinlinemathZ
\lthtmlcheckvsize\clearpage}

{\newpage\clearpage
\lthtmlinlinemathA{tex2html_wrap_inline2788}%
$ \sim$%
\lthtmlinlinemathZ
\lthtmlcheckvsize\clearpage}

{\newpage\clearpage
\lthtmlinlinemathA{tex2html_wrap_indisplay2790}%
$\displaystyle (A,B,C) \sim (A',B',C') \iff \{A,B,C\} = \{A',B',C' \}$%
\lthtmlindisplaymathZ
\lthtmlcheckvsize\clearpage}

{\newpage\clearpage
\lthtmlinlinemathA{tex2html_wrap_inline2792}%
$ A,B,C \in \mathbb{R}^2$%
\lthtmlinlinemathZ
\lthtmlcheckvsize\clearpage}

{\newpage\clearpage
\lthtmldisplayA{displaymath2794}%
\begin{displaymath}\begin{array}[t]{rcrcl}
                 f & : & \mathcal{T}  & \to     & (\mathbb{R}^2 \times \mathbb{R}^2 \times \mathbb{R}^2)/_{\sim} \\
                   &   & T            & \mapsto & \widetilde{T}
          \end{array}\end{displaymath}%
\lthtmldisplayZ
\lthtmlcheckvsize\clearpage}

{\newpage\clearpage
\lthtmlinlinemathA{tex2html_wrap_inline2796}%
$ \widetilde{T}$%
\lthtmlinlinemathZ
\lthtmlcheckvsize\clearpage}

{\newpage\clearpage
\lthtmlinlinemathA{tex2html_wrap_inline2798}%
$ T$%
\lthtmlinlinemathZ
\lthtmlcheckvsize\clearpage}

{\newpage\clearpage
\lthtmlinlinemathA{tex2html_wrap_inline2802}%
$ \overline{\overline{\mathcal{T}}} \leqslant
\overline{\overline{(\mathbb{R}^2 \times \mathbb{R}^2 \times \mathbb{R}^2)/_{\sim}}}
\leqslant \overline{\overline{\mathbb{R}^2 \times \mathbb{R}^2 \times \mathbb{R}^2}} = c$%
\lthtmlinlinemathZ
\lthtmlcheckvsize\clearpage}

{\newpage\clearpage
\lthtmlinlinemathA{tex2html_wrap_inline2804}%
$ \overline{\overline{\mathcal{T}}} = c$%
\lthtmlinlinemathZ
\lthtmlcheckvsize\clearpage}

{\newpage\clearpage
\lthtmldisplayA{displaymath2806}%
\begin{displaymath}\begin{array}[t]{rcrcl}
                         f & : & \mathcal{D}  & \to     & \mathbb{R} \times \mathbb{R} \times \mathbb{R} \\
                     &   & d                    & \mapsto & \left\lbrace \begin{array}[]{ll}
                                                                                    (m,b,0) & \textrm{si $d$ s'�crit sous la forme $y=mx+b$} \\
                                                                                    (0,0,r) & \textrm{si $d$ s'�crit sous la forme $x=r$}
                                                                                    \end{array}
                                                                                    \right.
                   \end{array}\end{displaymath}%
\lthtmldisplayZ
\lthtmlcheckvsize\clearpage}

{\newpage\clearpage
\lthtmlinlinemathA{tex2html_wrap_inline2810}%
$ \overline{\overline{\mathcal{D}}} \leqslant \overline{\overline{\mathbb{R} \times \mathbb{R} \times \mathbb{R}}} = c$%
\lthtmlinlinemathZ
\lthtmlcheckvsize\clearpage}

{\newpage\clearpage
\lthtmldisplayA{displaymath2812}%
\begin{displaymath}\begin{array}[t]{rcrcl}
                    g & : & \mathbb{R}  & \to     & \mathcal{D} \\
                  &   & b                 & \mapsto & d_b = \{ (x,y) \in \mathbb{R}^2 \mid y = x + b \}
               \end{array}\end{displaymath}%
\lthtmldisplayZ
\lthtmlcheckvsize\clearpage}

{\newpage\clearpage
\lthtmlinlinemathA{tex2html_wrap_inline2816}%
$ \overline{\overline{\mathbb{R}}} \leqslant \overline{\overline{\mathcal{D}}}$%
\lthtmlinlinemathZ
\lthtmlcheckvsize\clearpage}

{\newpage\clearpage
\lthtmlinlinemathA{tex2html_wrap_inline2818}%
$ c \leqslant \overline{\overline{\mathcal{D}}}$%
\lthtmlinlinemathZ
\lthtmlcheckvsize\clearpage}

{\newpage\clearpage
\lthtmlinlinemathA{tex2html_wrap_inline2820}%
$ \overline{\overline{\mathcal{D}}} = c$%
\lthtmlinlinemathZ
\lthtmlcheckvsize\clearpage}

{\newpage\clearpage
\lthtmlinlinemathA{tex2html_wrap_inline2824}%
$ \mathbb{R}^2 = \bigcup\limits_{i \in I} d_i$%
\lthtmlinlinemathZ
\lthtmlcheckvsize\clearpage}

{\newpage\clearpage
\lthtmlinlinemathA{tex2html_wrap_inline2826}%
$ d_i$%
\lthtmlinlinemathZ
\lthtmlcheckvsize\clearpage}

{\newpage\clearpage
\lthtmlinlinemathA{tex2html_wrap_inline2830}%
$ A = \left\{ (x,e^x) \in \mathbb{R}^2 \mid x \textrm{ est un nombre irrationnel} \right\}$%
\lthtmlinlinemathZ
\lthtmlcheckvsize\clearpage}

{\newpage\clearpage
\lthtmlinlinemathA{tex2html_wrap_inline2832}%
$ \overline{\overline{A}} = \overline{\overline{\mathcal{I}}} = c$%
\lthtmlinlinemathZ
\lthtmlcheckvsize\clearpage}

{\newpage\clearpage
\lthtmlinlinemathA{tex2html_wrap_inline2836}%
$ A$%
\lthtmlinlinemathZ
\lthtmlcheckvsize\clearpage}

{\newpage\clearpage
\lthtmlinlinemathA{tex2html_wrap_inline2846}%
$ \frac{\overline{\overline{A}}}{2}$%
\lthtmlinlinemathZ
\lthtmlcheckvsize\clearpage}

{\newpage\clearpage
\lthtmlinlinemathA{tex2html_wrap_inline2850}%
$ \overline{\overline{I}} \geqslant
\frac{\overline{\overline{A}}}{2} = c/2 = c$%
\lthtmlinlinemathZ
\lthtmlcheckvsize\clearpage}

{\newpage\clearpage
\lthtmlinlinemathA{tex2html_wrap_inline2854}%
$ \mathbb{R}^2$%
\lthtmlinlinemathZ
\lthtmlcheckvsize\clearpage}

{\newpage\clearpage
\lthtmlinlinemathA{tex2html_wrap_inline2856}%
$ D_{k,k'} = \{ (x,y) \in \mathbb{R}^2 \mid (x-k)^2 + (y-k')^2 \leqslant 4 \}$%
\lthtmlinlinemathZ
\lthtmlcheckvsize\clearpage}

{\newpage\clearpage
\lthtmlinlinemathA{tex2html_wrap_inline2858}%
$ \mathbb{R}^2 = \bigcup\limits_{k,k' \in \mathbb{Z}} D_{k,k'}$%
\lthtmlinlinemathZ
\lthtmlcheckvsize\clearpage}

{\newpage\clearpage
\lthtmlinlinemathA{tex2html_wrap_inline2860}%
$ f^{-1}(n) \subseteq \mathbb{R}$%
\lthtmlinlinemathZ
\lthtmlcheckvsize\clearpage}

{\newpage\clearpage
\lthtmlinlinemathA{tex2html_wrap_inline2862}%
$ \alpha = \overline{\overline{f^{-1}(n)}} \leqslant \overline{\overline{\mathbb{R}}} = c$%
\lthtmlinlinemathZ
\lthtmlcheckvsize\clearpage}

{\newpage\clearpage
\lthtmlinlinemathA{tex2html_wrap_inline2864}%
$ \alpha > \aleph_0$%
\lthtmlinlinemathZ
\lthtmlcheckvsize\clearpage}

{\newpage\clearpage
\lthtmlinlinemathA{tex2html_wrap_inline2866}%
$ \alpha \leqslant \aleph_0$%
\lthtmlinlinemathZ
\lthtmlcheckvsize\clearpage}

{\newpage\clearpage
\lthtmlinlinemathA{tex2html_wrap_inline2868}%
$ \mathbb{R} = \bigcup\limits_{n \in \mathbb{N}} f^{-1}(n)$%
\lthtmlinlinemathZ
\lthtmlcheckvsize\clearpage}

{\newpage\clearpage
\lthtmlinlinemathA{tex2html_wrap_inline2872}%
$ V= \{ n \in \mathbb{N} \mid n \notin f(n) \}$%
\lthtmlinlinemathZ
\lthtmlcheckvsize\clearpage}

{\newpage\clearpage
\lthtmlinlinemathA{tex2html_wrap_inline2886}%
$ \mathcal{D}$%
\lthtmlinlinemathZ
\lthtmlcheckvsize\clearpage}

{\newpage\clearpage
\lthtmlinlinemathA{tex2html_wrap_inline2888}%
$ c$%
\lthtmlinlinemathZ
\lthtmlcheckvsize\clearpage}

{\newpage\clearpage
\lthtmlinlinemathA{tex2html_wrap_inline2890}%
$ \overline{\overline{D \times D}} = c\cdot c = c$%
\lthtmlinlinemathZ
\lthtmlcheckvsize\clearpage}

{\newpage\clearpage
\lthtmlinlinemathA{tex2html_wrap_inline2892}%
$ {\rm Im}(\sin) = \left[-1,1\right]$%
\lthtmlinlinemathZ
\lthtmlcheckvsize\clearpage}

{\newpage\clearpage
\lthtmlinlinemathA{tex2html_wrap_inline2894}%
$ \mathbb{R}$%
\lthtmlinlinemathZ
\lthtmlcheckvsize\clearpage}

{\newpage\clearpage
\lthtmlinlinemathA{tex2html_wrap_inline2898}%
$ \overline{\overline{{\rm Im}(\sin)}} = c$%
\lthtmlinlinemathZ
\lthtmlcheckvsize\clearpage}

{\newpage\clearpage
\lthtmlinlinemathA{tex2html_wrap_inline2900}%
$ \overline{\overline{{\cal F}}} = \overline{\overline{\mathbb{R}^{\mathbb{R}}}}
= \overline{\overline{\mathbb{R}}}^{\overline{\overline{\mathbb{R}}}} = c^c = (2^{\aleph_0})^c =
2^{\aleph_0 \cdot c} = 2^c$%
\lthtmlinlinemathZ
\lthtmlcheckvsize\clearpage}

{\newpage\clearpage
\lthtmlinlinemathA{tex2html_wrap_inline2902}%
$ \mathbb{N}_*$%
\lthtmlinlinemathZ
\lthtmlcheckvsize\clearpage}

{\newpage\clearpage
\lthtmlinlinemathA{tex2html_wrap_inline2906}%
$ \mathcal{S}$%
\lthtmlinlinemathZ
\lthtmlcheckvsize\clearpage}

{\newpage\clearpage
\lthtmlinlinemathA{tex2html_wrap_inline2908}%
$ \overline{\overline{\mathcal{S}}} = \overline{\overline{\mathbb{R}^{\mathbb{N}_*}}} = \overline{\overline{\mathbb{R}}}^{\overline{\overline{\mathbb{N}_*}}}
= c^{\aleph_0} = (2^{\aleph_0})^{\aleph_0} = 2^{\aleph_0 \cdot \aleph_0} = 2^{\aleph_0} = c$%
\lthtmlinlinemathZ
\lthtmlcheckvsize\clearpage}

{\newpage\clearpage
\lthtmlinlinemathA{tex2html_wrap_inline2910}%
$ \mathcal{S}_c$%
\lthtmlinlinemathZ
\lthtmlcheckvsize\clearpage}

{\newpage\clearpage
\lthtmlinlinemathA{tex2html_wrap_inline2912}%
$ \mathcal{S}_c \subseteq \mathcal{S}$%
\lthtmlinlinemathZ
\lthtmlcheckvsize\clearpage}

{\newpage\clearpage
\lthtmlinlinemathA{tex2html_wrap_inline2914}%
$ \overline{\overline{\mathcal{S}_c}} \leqslant \overline{\overline{\mathcal{S}}} = c$%
\lthtmlinlinemathZ
\lthtmlcheckvsize\clearpage}

{\newpage\clearpage
\lthtmldisplayA{displaymath2916}%
\begin{displaymath}\begin{array}[t]{rcrcl}
         f & : & \mathbb{R}  & \to     & \mathcal{S}_c \\
           &   & x                 & \mapsto & (x,x,x,x,x,\ldots )
          \end{array}\end{displaymath}%
\lthtmldisplayZ
\lthtmlcheckvsize\clearpage}

{\newpage\clearpage
\lthtmlinlinemathA{tex2html_wrap_inline2920}%
$ c = \overline{\overline{\mathbb{R}}} \leqslant \overline{\overline{\mathcal{S}_c}}$%
\lthtmlinlinemathZ
\lthtmlcheckvsize\clearpage}

{\newpage\clearpage
\lthtmlinlinemathA{tex2html_wrap_inline2922}%
$ \overline{\overline{\mathcal{S}_c}} = c$%
\lthtmlinlinemathZ
\lthtmlcheckvsize\clearpage}

{\newpage\clearpage
\lthtmldisplayA{displaymath2924}%
\begin{displaymath}\begin{array}[t]{rcrcl}
          h & : & \mathcal{C}   & \to     & \mathbb{R}^{\mathbb{Q}} \\
            &   & f                       & \mapsto & \begin{array}[t]{rcrcl}
                                                s & : & \mathbb{Q} & \to     & \mathbb{R} \\
                                                &   &  x             & \mapsto & f(x).
                                            \end{array}
          \end{array}\end{displaymath}%
\lthtmldisplayZ
\lthtmlcheckvsize\clearpage}

{\newpage\clearpage
\lthtmlinlinemathA{tex2html_wrap_inline2928}%
$ f,g \in \mathcal{C}$%
\lthtmlinlinemathZ
\lthtmlcheckvsize\clearpage}

{\newpage\clearpage
\lthtmlinlinemathA{tex2html_wrap_inline2930}%
$ h(f) = h(g)$%
\lthtmlinlinemathZ
\lthtmlcheckvsize\clearpage}

{\newpage\clearpage
\lthtmlinlinemathA{tex2html_wrap_inline2932}%
$ h(f)(x) = h(g)(x)$%
\lthtmlinlinemathZ
\lthtmlcheckvsize\clearpage}

{\newpage\clearpage
\lthtmlinlinemathA{tex2html_wrap_inline2934}%
$ x \in \mathbb{Q}$%
\lthtmlinlinemathZ
\lthtmlcheckvsize\clearpage}

{\newpage\clearpage
\lthtmlinlinemathA{tex2html_wrap_inline2936}%
$ f(x) = g(x)$%
\lthtmlinlinemathZ
\lthtmlcheckvsize\clearpage}

{\newpage\clearpage
\lthtmlinlinemathA{tex2html_wrap_inline2940}%
$ f=g$%
\lthtmlinlinemathZ
\lthtmlcheckvsize\clearpage}

{\newpage\clearpage
\lthtmlinlinemathA{tex2html_wrap_inline2944}%
$ \overline{\overline{\mathcal{C}}} \leqslant \overline{\overline{\mathbb{R}^{\mathbb{Q}}}} =
\overline{\overline{\mathbb{R}}}^{\overline{\overline{\mathbb{Q}}}} = c^{\aleph_0} = c$%
\lthtmlinlinemathZ
\lthtmlcheckvsize\clearpage}

{\newpage\clearpage
\lthtmldisplayA{displaymath2946}%
\begin{displaymath}\begin{array}[t]{rcrcl}
          f & : & \mathbb{R}    & \to     & \mathcal{C} \\
            &   & x_0                   & \mapsto & \begin{array}[t]{rcrcl}
                                            f_{x_0} & : & \mathbb{R} & \to     & \mathbb{R} \\
                                                    &   &  x             & \mapsto & x_0
                                            \end{array}
          \end{array}\end{displaymath}%
\lthtmldisplayZ
\lthtmlcheckvsize\clearpage}

{\newpage\clearpage
\lthtmlinlinemathA{tex2html_wrap_inline2948}%
$ c = \overline{\overline{\mathbb{R}}} \leqslant \overline{\overline{\mathcal{C}}}$%
\lthtmlinlinemathZ
\lthtmlcheckvsize\clearpage}

{\newpage\clearpage
\lthtmlinlinemathA{tex2html_wrap_inline2950}%
$ \overline{\overline{\mathcal{C}}} = c$%
\lthtmlinlinemathZ
\lthtmlcheckvsize\clearpage}

{\newpage\clearpage
\lthtmlinlinemathA{tex2html_wrap_inline2952}%
$ \mathcal{M}_n$%
\lthtmlinlinemathZ
\lthtmlcheckvsize\clearpage}

{\newpage\clearpage
\lthtmlinlinemathA{tex2html_wrap_inline2954}%
$ \{a,b \}$%
\lthtmlinlinemathZ
\lthtmlcheckvsize\clearpage}

{\newpage\clearpage
\lthtmlinlinemathA{tex2html_wrap_inline2958}%
$ \overline{\overline{\mathcal{M}_n}} = 2^n$%
\lthtmlinlinemathZ
\lthtmlcheckvsize\clearpage}

{\newpage\clearpage
\lthtmlinlinemathA{tex2html_wrap_inline2960}%
$ \mathcal{M} = \bigcup\limits_{n \in \mathbb{N}}
\mathcal{M}_n$%
\lthtmlinlinemathZ
\lthtmlcheckvsize\clearpage}

{\newpage\clearpage
\lthtmldisplayA{displaymath2962}%
\begin{displaymath}\begin{array}[t]{rcrcl}
                                                        g & : & \mathbb{N}_*  & \to     & \mathcal{M} \\
                                                    &   & n                     & \mapsto & \underbrace{aaa\ldots a}_{\textrm{$n$\  fois}}
                                                    \end{array}\end{displaymath}%
\lthtmldisplayZ
\lthtmlcheckvsize\clearpage}

{\newpage\clearpage
\lthtmlinlinemathA{tex2html_wrap_inline2964}%
$ \overline{\overline{\mathbb{N}}} \leqslant \overline{\overline{\mathcal{M}}}$%
\lthtmlinlinemathZ
\lthtmlcheckvsize\clearpage}

{\newpage\clearpage
\lthtmlinlinemathA{tex2html_wrap_inline2966}%
$ \overline{\overline{\mathcal{M}}} = \aleph_0$%
\lthtmlinlinemathZ
\lthtmlcheckvsize\clearpage}

{\newpage\clearpage
\lthtmlinlinemathA{tex2html_wrap_inline2968}%
$ {\cal D}$%
\lthtmlinlinemathZ
\lthtmlcheckvsize\clearpage}

{\newpage\clearpage
\lthtmlinlinemathA{tex2html_wrap_inline2970}%
$ \left]0,1\right]$%
\lthtmlinlinemathZ
\lthtmlcheckvsize\clearpage}

{\newpage\clearpage
\lthtmlinlinemathA{tex2html_wrap_inline2974}%
$ \overline{\overline{\mathcal{D}}} \leqslant c$%
\lthtmlinlinemathZ
\lthtmlcheckvsize\clearpage}

{\newpage\clearpage
\lthtmlinlinemathA{tex2html_wrap_inline2976}%
$ x$%
\lthtmlinlinemathZ
\lthtmlcheckvsize\clearpage}

{\newpage\clearpage
\lthtmlinlinemathA{tex2html_wrap_inline2980}%
$ x=0,x_1x_2\ldots x_n\ldots$%
\lthtmlinlinemathZ
\lthtmlcheckvsize\clearpage}

{\newpage\clearpage
\lthtmldisplayA{displaymath2982}%
\begin{displaymath}\begin{array}[t]{rcrcl}
          f & : & \mathcal{D}                           & \to     & \{0,1,2,4,5,6,7,8,9\}^{\mathbb{N}} \\
            &   & 0,x_1x_2\ldots x_n\ldots  & \mapsto & \begin{array}[t]{rcrcl}
                                                                        h & : & \mathbb{N} & \to     & \{0,1,2,4,5,6,7,8,9\} \\
                                                                        &   &  m             & \mapsto & x_m.
                                         \end{array}
          \end{array}\end{displaymath}%
\lthtmldisplayZ
\lthtmlcheckvsize\clearpage}

{\newpage\clearpage
\lthtmlinlinemathA{tex2html_wrap_inline2986}%
$ \overline{\overline{\mathcal{D}}} \leqslant \overline{\overline{\{0,1,2,4,5,6,7,8,9\}^{\mathbb{N}}}}
= 9^{\aleph_0} \leqslant 16^{\aleph_0} = (2^4)^{\aleph_0} = 2^{4 \cdot \aleph_0} = 2^{\aleph_0} = c$%
\lthtmlinlinemathZ
\lthtmlcheckvsize\clearpage}

{\newpage\clearpage
\lthtmlinlinemathA{tex2html_wrap_inline2992}%
$ \mathbb{N} \times \mathbb{N}$%
\lthtmlinlinemathZ
\lthtmlcheckvsize\clearpage}

{\newpage\clearpage
\lthtmlinlinemathA{tex2html_wrap_inline2994}%
$ {\cal E} \subseteq \mathcal{P}(\mathbb{N} \times \mathbb{N})$%
\lthtmlinlinemathZ
\lthtmlcheckvsize\clearpage}

{\newpage\clearpage
\lthtmlinlinemathA{tex2html_wrap_inline2996}%
$ \overline{\overline{\mathcal{E}}} \leqslant \overline{\overline{\mathcal{P}(\mathbb{N} \times \mathbb{N})}} = 2^{\overline{\overline{\mathbb{N} \times \mathbb{N}}}}
= 2^{\aleph_0} = c$%
\lthtmlinlinemathZ
\lthtmlcheckvsize\clearpage}

{\newpage\clearpage
\lthtmlinlinemathA{tex2html_wrap_inline2998}%
$ A \in \mathcal{P}(\mathbb{N})$%
\lthtmlinlinemathZ
\lthtmlcheckvsize\clearpage}

{\newpage\clearpage
\lthtmlinlinemathA{tex2html_wrap_inline3002}%
$ R_A$%
\lthtmlinlinemathZ
\lthtmlcheckvsize\clearpage}

{\newpage\clearpage
\lthtmlinlinemathA{tex2html_wrap_inline3004}%
$ m R_A n$%
\lthtmlinlinemathZ
\lthtmlcheckvsize\clearpage}

{\newpage\clearpage
\lthtmlinlinemathA{tex2html_wrap_inline3006}%
$ m=n$%
\lthtmlinlinemathZ
\lthtmlcheckvsize\clearpage}

{\newpage\clearpage
\lthtmlinlinemathA{tex2html_wrap_inline3008}%
$ m$%
\lthtmlinlinemathZ
\lthtmlcheckvsize\clearpage}

{\newpage\clearpage
\lthtmldisplayA{displaymath3016}%
\begin{displaymath}\begin{array}[t]{rcrcl}
                 g & : & \mathcal{P}(\mathbb{N}) & \to     & \mathcal{E} \\
                     &   & A                               & \mapsto & R_A
                    \end{array}\end{displaymath}%
\lthtmldisplayZ
\lthtmlcheckvsize\clearpage}

{\newpage\clearpage
\lthtmlinlinemathA{tex2html_wrap_inline3020}%
$ \overline{\overline{\mathcal{P}(\mathbb{N})}} \leqslant \overline{\overline{\mathcal{E}}}$%
\lthtmlinlinemathZ
\lthtmlcheckvsize\clearpage}

{\newpage\clearpage
\lthtmlinlinemathA{tex2html_wrap_inline3022}%
$ \overline{\overline{\mathcal{E}}} \geqslant \overline{\overline{\mathcal{P}(\mathbb{N})}} = 2^{\overline{\overline{\mathbb{N}}}} = 2^{\aleph_0} = c$%
\lthtmlinlinemathZ
\lthtmlcheckvsize\clearpage}

{\newpage\clearpage
\lthtmlinlinemathA{tex2html_wrap_inline3024}%
$ \overline{\overline{\mathcal{E}}} = c$%
\lthtmlinlinemathZ
\lthtmlcheckvsize\clearpage}

{\newpage\clearpage
\lthtmldisplayA{displaymath3026}%
\begin{displaymath}\begin{array}[t]{rcrcl}
      h & : & \{0,1,2\}^{\mathbb{R}}  & \to     & \Pi_3 \\
        &   & f                       & \mapsto & \left\{ A,B,C \right\}
      \end{array}\end{displaymath}%
\lthtmldisplayZ
\lthtmlcheckvsize\clearpage}

{\newpage\clearpage
\lthtmldisplayA{displaymath3028}%
\begin{displaymath}\begin{array}[t]{l}
    A = \{x \in \mathbb{R} \mid f(x)=0\} \\
    B = \{x \in \mathbb{R} \mid f(x)=1\} \\
    C = \{x \in \mathbb{R} \mid f(x)=2\}
    \end{array}\end{displaymath}%
\lthtmldisplayZ
\lthtmlcheckvsize\clearpage}

{\newpage\clearpage
\lthtmlinlinemathA{tex2html_wrap_inline3032}%
$ \overline{\overline{\Pi_3}} \leqslant \overline{\overline{\{0,1,2\}^{\mathbb{R}}}} =
3^c \leqslant 4^c = (2^2)^c = 2^{2c} = 2^c$%
\lthtmlinlinemathZ
\lthtmlcheckvsize\clearpage}

{\newpage\clearpage
\lthtmlinlinemathA{tex2html_wrap_inline3034}%
$ A \in \mathcal{P}(\mathbb{R}^+_*)$%
\lthtmlinlinemathZ
\lthtmlcheckvsize\clearpage}

{\newpage\clearpage
\lthtmlinlinemathA{tex2html_wrap_inline3036}%
$ \left\{ A, \{0\}, \mathbb{R}\setminus (A \cup \{0\}) \right\}$%
\lthtmlinlinemathZ
\lthtmlcheckvsize\clearpage}

{\newpage\clearpage
\lthtmlinlinemathA{tex2html_wrap_inline3040}%
$ \left\{ \left\{ A, \{0\}, \mathbb{R}\setminus (A \cup \{0\}) \right\} \mid A \in \mathcal{P}(\mathbb{R}^+_*) \right\}
\subseteq \Pi_3$%
\lthtmlinlinemathZ
\lthtmlcheckvsize\clearpage}

{\newpage\clearpage
\lthtmlinlinemathA{tex2html_wrap_inline3042}%
$ \overline{\overline{\left\{ \left\{ A, \{0\}, \mathbb{R}\setminus (A \cup \{0\}) \right\} \mid A \in \mathcal{P}(\mathbb{R}^+_*) \right\}}}
\leqslant \overline{\overline{\Pi_3}}$%
\lthtmlinlinemathZ
\lthtmlcheckvsize\clearpage}

{\newpage\clearpage
\lthtmlinlinemathA{tex2html_wrap_inline3044}%
$ \overline{\overline{\mathcal{P}(\mathbb{R}^+_*)}} \leqslant \overline{\overline{\Pi_3}}$%
\lthtmlinlinemathZ
\lthtmlcheckvsize\clearpage}

{\newpage\clearpage
\lthtmlinlinemathA{tex2html_wrap_inline3046}%
$ 2^c \leqslant \overline{\overline{\Pi_3}}$%
\lthtmlinlinemathZ
\lthtmlcheckvsize\clearpage}

{\newpage\clearpage
\lthtmlinlinemathA{tex2html_wrap_inline3048}%
$ \overline{\overline{\Pi_3}} = 2^c$%
\lthtmlinlinemathZ
\lthtmlcheckvsize\clearpage}


\end{document}
